\section{Discussion and Limitations}
\label{sec:discussion}

\paragraph{Implications for pipeline authors.}
Three practices follow from the data, each at the scale of this corpus.
First, remove fences before validating, then revalidate and record what
remains. Fence removal is a partial repair: it recovers the
pure-formatting share of wrapped failures, which ranged from most of
Sonnet~5's to under half of Haiku~4.5's, and a pipeline that counts it as
a fix would report the compound remainder as solved. Second, when
diagnostic clarity matters, represent identifier length separately. A
\code{\{2,30\}} quantifier turns an over-long identifier into a pattern
violation, and the natural remedy for a pattern violation, loosening the
character class, leaves it untouched; a separate \code{max\_length}
separates the causes at validation time. Third, measure a field's
observed length distribution before fixing its cap. Sonnet~4.6's
dominant mode came from caps sitting modestly below its output lengths,
and Haiku~4.5's incomplete runs all ended on length. Such a cap may still
be a real downstream requirement; the measurement identifies the
mismatch the application must resolve, not that the cap is wrong.

\paragraph{Durability.}
A reasonable objection is that providers are engineering these modes
away. Our cross-generation data are descriptive and do not isolate model
generation from other configuration differences, including collection
date, sampling parameters and endpoint. With that caveat, they point in
both directions: \cat{schema\_mirroring} dominates gpt-4o and was not
observed in the GPT-5.6 configurations we tested, while
\cat{markdown\_wrap} was not observed in Sonnet~4.6 and dominates
Sonnet~5 and Haiku~4.5. What differs across these models is which
category dominates. A vocabulary for the categories remains useful as the
distribution moves, which is a weaker claim than predicting that any
particular mode persists.

\paragraph{Scope and generality.}
The schemas are ours, and one ablation showed that a cap we set was
binding below the output lengths of one model; a different author would
produce a different distribution over categories, and the domain
comparison shows that a different task would too. We cover two vendors
by design, trading breadth for several generations and capability tiers
per vendor, which the within-family comparisons above require. A third
vendor may exhibit modes our taxonomy lacks. The taxonomy is an operational
diagnostic for this paradigm, not a complete account of structured-output
failure, and most failures in this corpus come from a single model. Questions are fixed dataset
prefixes, and the design is not fully crossed. Most endpoints returned an
alias rather than a pinned version, so later calls to the same identifier
may not reproduce these distributions. We evaluate conformance only, not
whether extracted content is correct.

\paragraph{Measurement.}
Repair attempts, which are conditioned on a prior failure, make up a
model-dependent share of each attempt rate; first-attempt rates and run
completion are reported separately for that reason (\cref{app:breakdowns}).
Five models reject an explicit temperature, and the Anthropic and OpenAI
branches use different token ceilings. Some questions were run more than
once. Keeping only the earliest run per question changes any cell's
failure rate by at most \pn{robust.within_cell.max_abs_delta_pp}
percentage points (\cref{app:manifest}); the Haiku~4.5 FinanceBench cell
rests on \pn{c.claude-haiku-4-5@medium@financebench.questions} distinct
questions. Classifier agreement is development-set agreement from a
single labeller who wrote the classifier. Primary categories are
single-valued, and we quantify co-occurrence only beneath fences.

\paragraph{What is not tested.}
Among attempts whose model-level validators ran, no cross-field
violation occurred (\cref{app:crossfield}), but attempts rejected at an
earlier stage, including during JSON parsing, never reached those checks,
so the zero describes a filtered subset of the node-level validators we
implemented. Cross-node referential integrity---whether a later node's
identifiers appear in an earlier node's entities---is not implemented in
our schemas. It is expressible in Pydantic by passing the earlier output
as validation context (\code{model\_validate\_json(data, context=\ldots)}),
and is the natural next measurement.
